\documentclass{article}
\usepackage{amsmath,amssymb}
\usepackage{xurl}
\usepackage[hidelinks]{hyperref}
% Shared commands for notes in docs/paper-gaps/.

\DeclareUnicodeCharacter{2080}{\ifmmode{}_{0}\else\textsubscript{0}\fi}
\DeclareUnicodeCharacter{2081}{\ifmmode{}_{1}\else\textsubscript{1}\fi}
\DeclareUnicodeCharacter{2082}{\ifmmode{}_{2}\else\textsubscript{2}\fi}
\DeclareUnicodeCharacter{2099}{\ifmmode{}_{n}\else\textsubscript{n}\fi}

\newcommand{\Lean}{\textsc{Lean}}
\ExplSyntaxOn
\NewDocumentCommand{\leanid}{m}{
  \ifmmode
    \textsf{\detokenize{#1}}
  \else
    \group_begin:
    \urlstyle{sf}
    \tl_set:Nn \l_tmpa_tl {#1}
    \tl_replace_all:Nnn \l_tmpa_tl {\_} {_}
    \exp_args:NV \path \l_tmpa_tl
    \group_end:
  \fi
}
\ExplSyntaxOff
\providecommand{\tr}{\operatorname{tr}}
\newcommand{\tnleanIssueBase}{https://github.com/LionSR/TNLean/issues/}
\newcommand{\tnleanPrBase}{https://github.com/LionSR/TNLean/pull/}
\newcommand{\tnleanDocBase}{https://sirui-lu.com/TNLean/docs/}
\newcommand{\ghissue}[1]{\href{\tnleanIssueBase#1}{GitHub issue \##1}}
\newcommand{\ghpr}[1]{\href{\tnleanPrBase#1}{PR \##1}}
\newcommand{\doclink}[1]{\href{\tnleanDocBase#1.html}{\path{#1}}}

% Dirac notation and the identity operator, as in blueprint/src/macros/common.tex.
% \providecommand keeps note-local definitions authoritative.
\usepackage{dsfont}
\providecommand{\ket}[1]{|#1\rangle}
\providecommand{\bra}[1]{\langle #1|}
\providecommand{\braket}[2]{\langle #1 | #2 \rangle}
\providecommand{\ketbra}[2]{|#1\rangle\!\langle #2|}
\providecommand{\Id}{\mathds{1}}
\providecommand{\one}{\mathds{1}}
\providecommand{\C}{\mathbb{C}}
\providecommand{\MN}[1]{M_{#1}(\C)}
\providecommand{\MD}{M_D(\C)}

% Verdict marker: \gapnote{<kind>}{<status>}. Kind is the policy
% classification (clarification, local-correction, scope-restriction,
% unfaithful, false-source, open-gap); status is open, wip (elimination
% underway), resolved, or historical. Machine-read by the site index;
% prints a verdict line.
\newcommand{\gapnote}[2]{\par\noindent\textbf{Verdict:} #1 (#2).\par}

\title{Peripheral Roots after Blocking a Periodic Tensor}
\author{}
\date{2026-09-30}
\begin{document}
\maketitle
\gapnote{clarification}{resolved}

\section{The printed expression}
The proof of Lemma~6 in~\cite{DeLasCuevas2017Irreducible}, at
\path{Papers/1708.00029/main.tex:452--456}, writes the peripheral spectrum
of $\mathcal E_A^p$ as $\{\omega^l:0\leq l<m/r\}$, where
$\omega=e^{2\pi i/m}$ and $r=\gcd(m,p)$. The exponent is missing the
factor $p$. Throughout this calculation $m,p>0$.

\section{Correction}
Put $q=m/r$. Spectral mapping gives
\begin{align}
    \operatorname{spec}_{\mathrm{per}}(\mathcal E_A^p)
     & =\{\lambda^p:\lambda^m=1\}
    =\{\omega^{pl}:0\leq l<q\}
    =\{z\in\mathbb C:z^q=1\}.
\end{align}
Indeed, $\omega^p$ has order $m/\gcd(m,p)$. For example, when
$m=4$ and $p=2$, the printed set is $\{1,i\}$ whereas the correct set
is $\{1,-1\}$. This is a correction to the displayed description of the
roots, not a change to the period $m/r$ asserted in the lemma.

\section{Formal statement}
\leanid{MPSTensor.peripheralEigenvalues_blockTensor} proves the spectral
mapping identity for every tensor and positive blocking length.
\leanid{MPSTensor.IsPeriodic.peripheral_blockTensor} identifies the resulting
roots for a periodic tensor. Both are in
\doclink{TNLean/MPS/Periodic/BlockingSpectrum.lean}.
The irreducibility of each compressed block is a separate assertion;
this note does not certify that remaining part of Lemma~6.

\bibliographystyle{alpha}
\bibliography{references}
\end{document}
